\section{Expected Contributions}

\begin{itemize}
  \item A clear formulation of \textbf{retrieval-conditioned reward modeling} for robot manipulation.
  \item A concrete simulator-derived supervision scheme for reward labels and reward-quality evaluation.
  \item An experimental protocol for measuring the effect of retrieval on both reward quality and downstream policy quality.
  \item A systematic analysis of \textbf{retrieval granularity} in long-horizon manipulation.
  \item A simulation-first baseline that can later support extensions to real robots, VLAs, or world models.
\end{itemize}

\section{Proposed Execution Roadmap}

\begin{longtable}{p{2.1cm}p{4.0cm}p{7.6cm}}
  \toprule
  Phase & Goal & Expected output \\
  \midrule
  P1 & Lock benchmark, OOD protocol, and supervision sources & Choose `LIBERO` as the primary benchmark, define ID/OOD splits, and map simulator predicates to stage-progress labels \\
  P2 & Build the memory bank & Retrieval dataset with task instructions, stage descriptions, clips/sub-trajectories, optional manual text, and a vector index \\
  P3 & Train the reward model & Retrieval-conditioned reward head trained with progress regression and auxiliary pairwise ranking \\
  P4 & Integrate with policy learning & Offline RL with IQL using learned reward, plus reward-quality diagnostics and downstream policy metrics \\
  P5 & Analyze and write up & Ablations over context quality, retrieval granularity, OOD robustness, compute cost, and future work \\
  \bottomrule
\end{longtable}

\section{Future Extensions}

Once the core contribution is stable, the project can expand along three directions:
\begin{enumerate}
  \item \textbf{Real robots / LeRobot / SO-101}: keep the present proposal as the method layer, with hardware used later only for deployability validation.
  \item \textbf{Connection to VLA}: use the reward model as an evaluation or post-training layer for backbones such as OpenVLA or OpenVLA-OFT \citep{kim2024openvla, wang2025openvlaoft}.
  \item \textbf{Connection to world models}: replace retrieval-conditioned reward on the current state with reward over imagined rollouts, moving closer to the VLA-RFT direction \citep{li2025vlarft}.
\end{enumerate}

\section{Conclusion}

The topic ``Retrieval-Augmented Reward Modeling for Long-Horizon Robot Manipulation in Simulation''
is \textbf{sufficiently novel if framed correctly}, \textbf{feasible if kept simulation-first},
and \textbf{well aligned with active research trends} in embodied retrieval,
demonstration-conditioned learning, and robotics reward modeling.

The key is not to claim a broad ``robotic RAG'' system,
but to show that multimodal retrieval can make reward denser,
more reliable, and more useful for long-horizon manipulation learning in simulation.
If that is demonstrated convincingly, the resulting system becomes a solid base
for later extensions toward VLAs, world models, or real robots.
